\documentclass[../main.tex]{subfiles}

\begin{document}
\section{Conclusion}

\subsection{Summary of Findings}
\noindent
On the PathMNIST test set the
\gls{PCA}-reduced \gls{RBF} \gls{SVM} (69.01\% accuracy,
68.26\% macro-$F_{1}$) and the convolutional network
(68.41\% / 64.69\%) finished within 0.60
percentage points of each other, while the multilayer perceptron lagged
notably (34.51\% / 29.47\%). The ranking confirms one
expectation and challenges another. As predicted, the \gls{MLP} pays a
real price for treating each pixel as an independent feature and ignoring
the spatial structure of \gls{HE} histology. Less expected is the near
parity between the kernel \gls{SVM} and the small \gls{CNN}: at
PathMNIST's 28$\times$28 resolution and with a deliberately small
two-block \gls{CNN}, an \gls{RBF} kernel applied on top of a
30-dimensional \gls{PCA} embedding captures most of the texture cues
that drive classification, leaving little room for the \gls{CNN} to
pull ahead. The \gls{CNN} retains a clear advantage on inference
latency ($0.36\,\text{s}$ versus $6.02\,\text{s}$ for
8{,}000 test images), so it remains the more sensible choice for a
deployment that has to score images quickly.

\subsection{Limitations}
\noindent
Several caveats should temper these conclusions. First, the search grids
were deliberately small to fit the available compute; broader grids over
network depth, weight decay and learning-rate schedules could change the
relative ranking, especially between the \gls{MLP} and the \gls{CNN}.
Second, the SVM was tuned on a stratified 5{,}000-sample subset to keep
the kernel computation tractable; the full-data refit improves accuracy
modestly but might benefit further from a re-tuned grid at the larger
scale, particularly for the kernel bandwidth $\gamma$. Third, we did not
apply data augmentation; on a larger study the marginal benefit of
augmentation versus the increased training-time budget should be measured
explicitly. Finally, the PathMNIST 28$\times$28 resolution is much
coarser than the source slides — extrapolating these conclusions to
clinical-resolution images would require re-running the comparison on a
higher-resolution variant such as PathMNIST-224.

\subsection{Future Work}
\noindent
Three concrete directions would strengthen the study. (i) Replace the
exhaustive grid for the neural networks with a Bayesian search over a
larger space (depth, weight decay, optimiser, schedule) to test whether
the present \gls{CNN} is on the Pareto frontier of accuracy versus
compute~\cite{bergstra2012random}. (ii) Re-evaluate the comparison on
PathMNIST-224, where convolutional inductive biases should matter more
and where the \gls{SVM} would need either a strong feature extractor
(e.g.~\gls{PCA} on a pre-trained backbone's features) or to be replaced
by a linear baseline, giving a sharper picture of the model
families' scaling behaviour. (iii) Calibrate per-class predictions and
report calibration metrics (e.g.~expected calibration error) — clinical
deployment of any of these models would require uncertainty estimates,
and the three families differ markedly in how their probability outputs
behave.

\end{document}
